% ---- begin paper/tables/t2rep_gillmore.tex
\begin{table}[htbp]\centering
\caption{Replication of Gillmore's (2026) Main Table}
\label{tab:gillmore}
\begin{threeparttable}
\begin{tabular}{@{}p{200pt}*{3}{>{\centering\arraybackslash}p{74pt}}@{}}
\toprule
 & (1) & (2) & (3) \\
\midrule
\panel{4}{Panel A. Receptive vocabulary: Peabody picture--word test (z)}
Affected $\times$ affected zone & $-$0.157\sym{**} & $-$0.158\sym{**} & $-$0.160\sym{**} \\
 & (0.074) & (0.074) & (0.072) \\
\addlinespace[3pt]
\hspace{1em}Observations & 14,183 & 14,183 & 14,183 \\
\addlinespace
\panel{4}{Panel B. Internalizing problems: caregiver CBCL checklist (z)}
Affected $\times$ affected zone & $-$0.091 & $-$0.092 & $-$0.103 \\
 & (0.080) & (0.080) & (0.082) \\
\addlinespace[3pt]
\hspace{1em}Observations & 11,628 & 11,628 & 11,628 \\
\midrule
Municipality fixed effects & Yes & Yes & Yes \\
Birth-cohort fixed effects & No & Yes & Yes \\
Household controls & No & No & Yes \\
\bottomrule
\end{tabular}
\begin{wpnotes}
\textit{Notes:} Replication on the public ELPI files of the medium-term table of Gillmore (2026), in his exact design: a single wave (2017), with affected children aged 7--11 compared with children conceived after the earthquake, aged 2--6. Column 1 includes municipality fixed effects only; column 2 adds birth-cohort fixed effects; column 3 the household controls. His published coefficients in column 3 are $-$0.174 for vocabulary and $-$0.129 for the CBCL. \zonedef{} Standard errors clustered by municipality in parentheses; evaluation weights. \starnote
\end{wpnotes}
\end{threeparttable}
\end{table}
% ---- end paper/tables/t2rep_gillmore.tex
